% Ready-to-paste supplementary results section (verified findings only). Requires booktabs; tables are in SUPPLEMENTARY_TABLES_GENERATED.tex.
% Figures: figures/FIG1..FIG7 (PDF). All numbers trace to SUPPLEMENTARY_EXPERIMENTS_PER_SEED.csv (run IDs per row).
\section{Supplementary Controlled Experiments}
\label{sec:supplementary}

All experiments in this section use the locked validation protocol (five paired seeds 12--16 unless stated; reserved test sets not accessed). We report paired mean differences with 95\% $t$-intervals (assuming approximately normal paired differences, which cannot be verified with five seeds) and the number of seeds in each direction. The exact sign-flip test cannot return a two-sided $p$ below $0.0625$ for $n=5$ ($0.25$ for $n=3$), and some intervals exclude zero while the exact test does not support $p<0.05$; since about 1{,}700 unadjusted comparisons were computed, an interval excluding zero is treated as exploratory evidence of direction only and no claim of statistical significance or superiority is made. CICIoT2023 sensitivity and scalability arms use three seeds (12--14).

\paragraph{Attack robustness.} Table~\ref{tab:supp-attack} and Fig.~\ref{fig:supp-attack} compare the MAD-plus-surrogate pipeline with FedAvg under static sign-flip (SF, 10/20/30\% exposure-controlled malicious clients), norm-preserving sign-flip (BSF, 20\%) and label-flip (LF, 20\%) attacks. FedAvg degrades with SF exposure (UNSW-NB15 selected-checkpoint ROC--AUC 0.982, 0.946, 0.887, 0.705 at 0, 10, 20, 30\%) and its last-round AUC falls to 0.50 at 30\%, whereas the proposed pipeline remains at 0.980--0.984 (UNSW-NB15) and 0.996 (CICIoT2023). At 30\% SF the paired difference is $+0.275$ (95\% interval $[+0.135,+0.415]$; higher in 5/5 seeds) on UNSW-NB15 and $+0.038$ ($[-0.014,+0.089]$; 5/5) on CICIoT2023. BSF and LF did not measurably harm either pipeline, and the filter retained 79--80\% of BSF updates, so these two attacks are weak in this setting rather than detected. Without an attack the two pipelines are indistinguishable and MAD rejects 5--7\% of honest updates. All attacks are static; \emph{adaptive} attackers were not evaluated and the results must not be read as robustness against them.

\paragraph{Sensitivity.} Table~\ref{tab:supp-sens} and Fig.~\ref{fig:supp-sens} vary the MAD coefficient, Dirichlet concentration, clipping threshold and $\chi_{\mathrm{rec}}$. The largest effect is data heterogeneity: Dirichlet $\alpha=0.1$ lowers selected-checkpoint AUC by $0.116$ ($[-0.211,-0.020]$; 5/5 seeds) on UNSW-NB15 and $0.066$ ($[-0.093,-0.040]$; 3/3) on CICIoT2023. Halving the clipping threshold lowers UNSW-NB15 AUC by $0.019$ (5/5 seeds). Other changes are below $0.012$ with intervals including zero. $\chi_{\mathrm{rec}}$ had no detectable effect (all $|\Delta|\le 0.0008$): under central Gaussian noise the effective weights are exactly uniform (Section~\ref{sec:aggregation-weights}), and in the non-private path the effect is below seed variation.

\paragraph{Routing and surrogates.} In the $2\times2$ design (Table~\ref{tab:supp-interaction}, Fig.~\ref{fig:supp-int}) the interaction contrast is $+0.0068$ ($[-0.0024,+0.0159]$) on UNSW-NB15 and $+0.0005$ ($[-0.0004,+0.0014]$) on CICIoT2023; no interaction is demonstrated. Surrogates carry 21--22\% of the aggregation weight under routing and 41--43\% under random selection (17\% of rounds contain only surrogates), without a demonstrated detection benefit.

\paragraph{Participation and fairness.} Table~\ref{tab:supp-participation} and Fig.~\ref{fig:supp-part} report selection Gini (0.73, 0.49, 0.21 for top-$k$, coverage-aware and uniform selection) and recall on never-trained-on client holdouts. Held-out recall of clients never selected equals that of selected clients (CICIoT2023 top-$k$: 0.905 and 0.905). More equal selection did not raise the worst eligible-client recall (UNSW-NB15 uniform versus top-$k$: $-0.138$, $[-0.260,-0.015]$, 5/5 seeds lower; CICIoT2023: $-0.006$, $[-0.088,+0.076]$) and increased surrogate-only rounds to 38--41\%. CICIoT2023 results use a leakage-safe client holdout and are not pooled with earlier CICIoT2023 runs.

\paragraph{Scalability and cost.} Table~\ref{tab:supp-scale} and Fig.~\ref{fig:supp-scale} (three seeds) show the lowest utility at $N=25,k=5$; increasing $N$ at $k=10$ did not improve utility beyond seed variation. Peak working set was constant (1.03~GB UNSW-NB15, 2.56~GB CICIoT2023); simulated update traffic scales linearly with $k$. Wall-clock values come from a shared CPU host and are indicative only.

\paragraph{$\chi_{\mathrm{rec}}=0$ and abstention.} When all retained updates have zero raw weight the round abstains (model and momentum unchanged, no noise, no release). This occurred in 1 round (seed 13) and 4 rounds (seed 15) of the $\chi_{\mathrm{rec}}=0$ DP runs and nowhere else; the remaining seeds are bit-identical to the reference. The diagnostic $\widehat{\varepsilon}$ is 9.86 and 9.43 (release-based) versus 9.9993 (attempted-round) for those seeds (Table~\ref{tab:supp-chi0}); both are RDP estimates and not certified differential-privacy guarantees.

% Auto-generated supplementary tables. Requires booktabs. Source: Supplementary_Experiments_Package CSVs (run IDs in PER_SEED.csv).

\begin{table*}[!t]
\caption{Attack robustness (five paired seeds 12--16; mean (SD); selected-checkpoint and last-round ROC--AUC). SF: sign-flip with exposure-controlled malicious fraction; BSF: norm-preserving sign-flip; LF: label flipping. All attacks are static and non-adaptive. SF 20\% rows are the reused locked Phase~6C arms.}
\label{tab:supp-attack}
\centering\tiny
\setlength{\tabcolsep}{2pt}
\begin{tabular}{@{}lllllllll@{}}
\toprule
Dataset & Condition & FedAvg AUC sel & MAD+surr AUC sel & FedAvg AUC last & MAD+surr AUC last & malicious retained & honest rejected & paired $\Delta$ AUC sel (ref $-$ FedAvg) [95\% CI] \\
\midrule
UNSW-NB15 & no attack & 0.9820 (0.0066) & 0.9829 (0.0037) & 0.9807 (0.0067) & 0.9811 (0.0066) & - & 0.068 (0.028) & +0.0009 [-0.0031, +0.0050] (n=5; 3$\downarrow$/2$\uparrow$) \\
UNSW-NB15 & SF 10\% & 0.9462 (0.0408) & 0.9807 (0.0069) & 0.8949 (0.0590) & 0.9800 (0.0069) & 0.085 (0.098) & 0.058 (0.034) & +0.0345 [-0.0097, +0.0787] (n=5; 1$\downarrow$/4$\uparrow$) \\
UNSW-NB15 & SF 20\% & 0.8870 (0.0856) & 0.9838 (0.0029) & 0.6043 (0.1430) & 0.9826 (0.0045) & 0.142 (0.075) & 0.053 (0.032) & +0.0968 [-0.0084, +0.2020] (n=5; 0$\downarrow$/5$\uparrow$) \\
UNSW-NB15 & SF 30\% & 0.7046 (0.1132) & 0.9799 (0.0026) & 0.5000 (0.0000) & 0.9777 (0.0036) & 0.221 (0.067) & 0.039 (0.037) & +0.2753 [+0.1353, +0.4152] (n=5; 0$\downarrow$/5$\uparrow$) \\
UNSW-NB15 & BSF 20\% & 0.9799 (0.0070) & 0.9800 (0.0057) & 0.9791 (0.0067) & 0.9770 (0.0051) & 0.803 (0.049) & 0.059 (0.028) & +0.0001 [-0.0074, +0.0076] (n=5; 2$\downarrow$/3$\uparrow$) \\
UNSW-NB15 & LF 20\% & 0.9778 (0.0088) & 0.9836 (0.0034) & 0.9776 (0.0067) & 0.9809 (0.0063) & 0.636 (0.110) & 0.053 (0.036) & +0.0058 [-0.0045, +0.0161] (n=5; 1$\downarrow$/4$\uparrow$) \\
CICIoT2023 & no attack & 0.9963 (0.0003) & 0.9962 (0.0004) & 0.9961 (0.0004) & 0.9961 (0.0004) & - & 0.053 (0.025) & -0.0001 [-0.0002, +0.0000] (n=5; 4$\downarrow$/1$\uparrow$) \\
CICIoT2023 & SF 10\% & 0.9949 (0.0010) & 0.9961 (0.0004) & 0.8071 (0.2463) & 0.9960 (0.0005) & 0.047 (0.032) & 0.048 (0.030) & +0.0012 [+0.0002, +0.0023] (n=5; 0$\downarrow$/5$\uparrow$) \\
CICIoT2023 & SF 20\% & 0.9813 (0.0211) & 0.9962 (0.0003) & 0.4625 (0.0843) & 0.9960 (0.0006) & 0.068 (0.062) & 0.032 (0.033) & +0.0149 [-0.0111, +0.0410] (n=5; 0$\downarrow$/5$\uparrow$) \\
CICIoT2023 & SF 30\% & 0.9583 (0.0420) & 0.9961 (0.0004) & 0.5000 (0.0000) & 0.9959 (0.0002) & 0.144 (0.067) & 0.035 (0.039) & +0.0377 [-0.0140, +0.0894] (n=5; 0$\downarrow$/5$\uparrow$) \\
CICIoT2023 & BSF 20\% & 0.9960 (0.0003) & 0.9961 (0.0004) & 0.9955 (0.0004) & 0.9949 (0.0019) & 0.787 (0.079) & 0.035 (0.036) & +0.0001 [-0.0001, +0.0003] (n=5; 1$\downarrow$/4$\uparrow$) \\
CICIoT2023 & LF 20\% & 0.9960 (0.0004) & 0.9961 (0.0003) & 0.6111 (0.3660) & 0.9958 (0.0008) & 0.406 (0.086) & 0.030 (0.034) & +0.0001 [-0.0001, +0.0002] (n=5; 1$\downarrow$/4$\uparrow$) \\
\bottomrule
\end{tabular}
\end{table*}

\begin{table*}[!t]
\caption{Sensitivity of selected-checkpoint ROC--AUC to MAD coefficient, Dirichlet concentration, clipping threshold and $\chi_{\mathrm{rec}}$ (paired difference to the parent arm, 95\% $t$-interval; $n$ = paired seeds; $\downarrow/\uparrow$ = seeds where the variant is lower/higher). With $n=5$ the smallest attainable two-sided exact sign-flip $p$ is 0.0625; no claim of statistical significance is made. $\chi_{\mathrm{rec}}$ is inert under the central-DP pipeline (see text).}
\label{tab:supp-sens}
\centering\tiny
\setlength{\tabcolsep}{2pt}
\begin{tabular}{@{}llll@{}}
\toprule
Dataset & Variant & Parent arm & paired $\Delta$ ROC-AUC (selected) [95\% CI] \\
\midrule
UNSW-NB15 & MAD k=1.5 & reference\_gaussian & -0.0030 [-0.0070, +0.0010] (n=5; 5$\downarrow$/0$\uparrow$) \\
UNSW-NB15 & MAD k=3.5 & reference\_gaussian & +0.0010 [-0.0002, +0.0022] (n=5; 0$\downarrow$/4$\uparrow$) \\
UNSW-NB15 & Dirichlet $\alpha$=0.1 & reference\_gaussian & -0.1159 [-0.2114, -0.0203] (n=5; 5$\downarrow$/0$\uparrow$) \\
UNSW-NB15 & Dirichlet $\alpha$=1.0 & reference\_gaussian & +0.0119 [-0.0051, +0.0288] (n=5; 0$\downarrow$/5$\uparrow$) \\
UNSW-NB15 & clip 0.5$\times$ & reference\_gaussian & -0.0187 [-0.0359, -0.0015] (n=5; 5$\downarrow$/0$\uparrow$) \\
UNSW-NB15 & clip 2$\times$ & reference\_gaussian & -0.0118 [-0.0289, +0.0052] (n=5; 4$\downarrow$/1$\uparrow$) \\
UNSW-NB15 & $\chi_{\mathrm{rec}}$=0 (DP) & reference\_gaussian & +0.0018 [-0.0018, +0.0054] (n=5; 0$\downarrow$/2$\uparrow$) \\
UNSW-NB15 & $\chi_{\mathrm{rec}}$=0 (no noise) & reference & -0.0008 [-0.0032, +0.0016] (n=5; 3$\downarrow$/2$\uparrow$) \\
UNSW-NB15 & $\chi_{\mathrm{rec}}$=0.25 (no noise) & reference & -0.0003 [-0.0010, +0.0004] (n=5; 4$\downarrow$/1$\uparrow$) \\
UNSW-NB15 & $\chi_{\mathrm{rec}}$=0.75 (no noise) & reference & -0.0000 [-0.0005, +0.0004] (n=5; 2$\downarrow$/3$\uparrow$) \\
CICIoT2023 & MAD k=1.5 & reference\_gaussian & -0.0003 [-0.0014, +0.0008] (n=3; 2$\downarrow$/1$\uparrow$) \\
CICIoT2023 & MAD k=3.5 & reference\_gaussian & +0.0000 [-0.0002, +0.0003] (n=3; 1$\downarrow$/1$\uparrow$) \\
CICIoT2023 & Dirichlet $\alpha$=0.1 & reference\_gaussian & -0.0664 [-0.0929, -0.0399] (n=3; 3$\downarrow$/0$\uparrow$) \\
CICIoT2023 & Dirichlet $\alpha$=1.0 & reference\_gaussian & +0.0001 [-0.0035, +0.0037] (n=3; 2$\downarrow$/1$\uparrow$) \\
CICIoT2023 & clip 0.5$\times$ & reference\_gaussian & -0.0005 [-0.0016, +0.0006] (n=3; 3$\downarrow$/0$\uparrow$) \\
CICIoT2023 & clip 2$\times$ & reference\_gaussian & -0.0005 [-0.0028, +0.0018] (n=3; 2$\downarrow$/1$\uparrow$) \\
CICIoT2023 & $\chi_{\mathrm{rec}}$=0 (DP) & reference\_gaussian & +0.0001 [-0.0002, +0.0004] (n=5; 1$\downarrow$/1$\uparrow$) \\
CICIoT2023 & $\chi_{\mathrm{rec}}$=0 (no noise) & reference & -0.0000 [-0.0001, +0.0001] (n=3; 2$\downarrow$/1$\uparrow$) \\
CICIoT2023 & $\chi_{\mathrm{rec}}$=0.25 (no noise) & reference & +0.0000 [-0.0000, +0.0001] (n=3; 0$\downarrow$/3$\uparrow$) \\
CICIoT2023 & $\chi_{\mathrm{rec}}$=0.75 (no noise) & reference & +0.0000 [-0.0001, +0.0001] (n=3; 1$\downarrow$/2$\uparrow$) \\
\bottomrule
\end{tabular}
\end{table*}

\begin{table*}[!t]
\caption{Routing$\times$surrogate $2\times2$ design (five seeds; matched traces within each routing state). Interaction = (surrogate effect | routing ON) $-$ (surrogate effect | routing OFF).}
\label{tab:supp-interaction}
\centering\tiny
\setlength{\tabcolsep}{2pt}
\begin{tabular}{@{}llllll@{}}
\toprule
Dataset & Cell & ROC-AUC sel mean (SD) & F1 sel mean (SD) & surrogate weight share & surrogate-only rounds (fraction) \\
\midrule
UNSW-NB15 & routing ON / surrogate OFF & 0.9822 (0.0056) & 0.9540 (0.0030) & 0.000 (0.000) & 0.000 (0.000) \\
UNSW-NB15 & routing ON / surrogate ON & 0.9838 (0.0029) & 0.9544 (0.0021) & 0.221 (0.048) & 0.033 (0.058) \\
UNSW-NB15 & routing OFF / surrogate OFF & 0.9737 (0.0208) & 0.9425 (0.0211) & 0.000 (0.000) & 0.000 (0.000) \\
UNSW-NB15 & routing OFF / surrogate ON & 0.9685 (0.0238) & 0.9436 (0.0144) & 0.434 (0.059) & 0.173 (0.083) \\
UNSW-NB15 & interaction contrast & +0.0068 [-0.0024, +0.0159] & exact sign-flip p=0.0625 (n=5, floor 0.0625) &  &  \\
CICIoT2023 & routing ON / surrogate OFF & 0.9962 (0.0003) & 0.9934 (0.0005) & 0.000 (0.000) & 0.000 (0.000) \\
CICIoT2023 & routing ON / surrogate ON & 0.9962 (0.0003) & 0.9934 (0.0002) & 0.208 (0.058) & 0.033 (0.058) \\
CICIoT2023 & routing OFF / surrogate OFF & 0.9960 (0.0003) & 0.9932 (0.0003) & 0.000 (0.000) & 0.000 (0.000) \\
CICIoT2023 & routing OFF / surrogate ON & 0.9955 (0.0009) & 0.9927 (0.0009) & 0.414 (0.074) & 0.173 (0.083) \\
CICIoT2023 & interaction contrast & +0.0005 [-0.0004, +0.0014] & exact sign-flip p=0.1250 (n=5, floor 0.0625) &  &  \\
\bottomrule
\end{tabular}
\end{table*}

\begin{table*}[!t]
\caption{Participation and held-out client recall (client holdouts never used for training; UNSW-NB15 from SUPP\_V1, CICIoT2023 from SUPP\_V2 with the leakage-safe holdout; the two protocols are not mixed). Recall is computed only for clients with at least 10 held-out attack rows (UNSW 43.8, CICIoT2023 48.4 of 50 on average).}
\label{tab:supp-participation}
\centering\tiny
\setlength{\tabcolsep}{2pt}
\begin{tabular}{@{}lllllllll@{}}
\toprule
Dataset & Policy / network & ROC-AUC sel & selection Gini & never-selected clients & worst eligible-client held-out recall & P10 recall & mean recall of never-selected & surrogate weight share \\
\midrule
UNSW-NB15 & top-k / indep. & 0.9787 (0.0084) & 0.731 (0.030) & 27.2 (2.0) & 0.916 (0.035) & 0.953 (0.015) & 0.970 (0.011) & 0.422 (0.063) \\
UNSW-NB15 & top-k / coupled & 0.9650 (0.0251) & 0.733 (0.022) & 27.0 (1.9) & 0.843 (0.131) & 0.897 (0.108) & 0.918 (0.098) & 0.508 (0.109) \\
UNSW-NB15 & coverage / indep. & 0.9649 (0.0199) & 0.486 (0.012) & 8.0 (2.9) & 0.861 (0.092) & 0.906 (0.090) & 0.935 (0.077) & 0.530 (0.073) \\
UNSW-NB15 & coverage / coupled & 0.9741 (0.0087) & 0.484 (0.011) & 8.4 (3.4) & 0.896 (0.058) & 0.933 (0.031) & 0.954 (0.025) & 0.556 (0.071) \\
UNSW-NB15 & uniform / indep. & 0.9426 (0.0320) & 0.205 (0.022) & 0.0 (0.0) & 0.779 (0.103) & 0.846 (0.108) & nan (nan) & 0.626 (0.138) \\
UNSW-NB15 & uniform / coupled & 0.9458 (0.0408) & 0.205 (0.022) & 0.0 (0.0) & 0.800 (0.159) & 0.851 (0.127) & nan (nan) & 0.630 (0.126) \\
CICIoT2023 & top-k / indep. & 0.9958 (0.0006) & 0.731 (0.030) & 27.2 (2.0) & 0.787 (0.048) & 0.886 (0.016) & 0.905 (0.011) & 0.404 (0.067) \\
CICIoT2023 & top-k / coupled & 0.9959 (0.0005) & 0.716 (0.025) & 26.2 (2.3) & 0.802 (0.048) & 0.894 (0.007) & 0.907 (0.007) & 0.433 (0.093) \\
CICIoT2023 & coverage / indep. & 0.9942 (0.0039) & 0.486 (0.012) & 8.0 (2.9) & 0.794 (0.019) & 0.877 (0.030) & 0.901 (0.030) & 0.485 (0.066) \\
CICIoT2023 & coverage / coupled & 0.9954 (0.0008) & 0.475 (0.025) & 6.4 (3.0) & 0.772 (0.065) & 0.885 (0.011) & 0.905 (0.013) & 0.506 (0.110) \\
CICIoT2023 & uniform / indep. & 0.9945 (0.0011) & 0.205 (0.022) & 0.0 (0.0) & 0.781 (0.067) & 0.862 (0.013) & nan (nan) & 0.583 (0.119) \\
CICIoT2023 & uniform / coupled & 0.9947 (0.0009) & 0.205 (0.022) & 0.0 (0.0) & 0.776 (0.076) & 0.874 (0.008) & nan (nan) & 0.611 (0.119) \\
\bottomrule
\end{tabular}
\end{table*}

\begin{table*}[!t]
\caption{Scalability of the proposed central-DP pipeline (seeds 12--14). Wall-clock and CPU seconds are measured on one shared CPU host and are not controlled benchmarks; bytes are simulated serialized-update accounting.}
\label{tab:supp-scale}
\centering\tiny
\setlength{\tabcolsep}{2pt}
\begin{tabular}{@{}lllllllll@{}}
\toprule
Dataset & Configuration (3 seeds 12-14) & ROC-AUC sel & wall s/run & CPU s/run & peak working set MB & local-train s/client-update & simulated MB on wire/run & surrogate weight share \\
\midrule
UNSW-NB15 & N=25,k=5 & 0.8675 (0.0210) & 62 & 192 & 1029 & 0.217 & 26.2 & 0.292 \\
UNSW-NB15 & N=25,k=10 & 0.9237 (0.0318) & 73 & 285 & 1029 & 0.159 & 52.4 & 0.313 \\
UNSW-NB15 & N=50,k=10 (n=3 seeds) & 0.9521 (0.0115) & 60 & 203 & 1029 & 0.103 & 52.4 & 0.274 \\
UNSW-NB15 & N=100,k=10 & 0.9205 (0.0200) & 47 & 120 & 1030 & 0.045 & 52.4 & 0.318 \\
UNSW-NB15 & N=100,k=20 & 0.9397 (0.0126) & 58 & 170 & 1029 & 0.043 & 104.7 & 0.248 \\
UNSW-NB15 & N=50,k=10 reference\_gaussian (n=5) & 0.9419 (0.0166) & 53 & 163 & 1030 & 0.089 & 52.4 & 0.306 \\
CICIoT2023 & N=25,k=5 & 0.9435 (0.0251) & 178 & 859 & 2557 & 1.025 & 14.3 & 0.286 \\
CICIoT2023 & N=25,k=10 & 0.9876 (0.0044) & 288 & 1570 & 2557 & 1.020 & 28.6 & 0.315 \\
CICIoT2023 & N=50,k=10 (n=3 seeds) & 0.9929 (0.0015) & 169 & 818 & 2556 & 0.443 & 28.6 & 0.271 \\
CICIoT2023 & N=100,k=10 & 0.9908 (0.0036) & 114 & 513 & 2557 & 0.229 & 28.6 & 0.317 \\
CICIoT2023 & N=100,k=20 & 0.9931 (0.0004) & 169 & 819 & 2556 & 0.227 & 57.2 & 0.220 \\
CICIoT2023 & N=50,k=10 reference\_gaussian (n=5) & 0.9930 (0.0014) & 180 & 843 & 2556 & 0.499 & 28.6 & 0.304 \\
\bottomrule
\end{tabular}
\end{table*}

\begin{table*}[!t]
\caption{$\chi_{\mathrm{rec}}=0$ under the central-DP pipeline with the no-valid-contribution rule (SUPP\_V2, five seeds). $\varepsilon$ values are diagnostic RDP estimates under a Poisson-subsampling assumption, not certified guarantees.}
\label{tab:supp-chi0}
\centering\tiny
\setlength{\tabcolsep}{2pt}
\begin{tabular}{@{}lllllllllllll@{}}
\toprule
Dataset & Seed & abstention rounds & round ids & noisy releases & attempted rounds & $\varepsilon$ release-based (diagnostic) & $\varepsilon$ attempted-round (diagnostic) & AUC sel $\chi$=0 & AUC sel reference\_gaussian & $\Delta$ sel & $\Delta$ last & V2 run id \\
\midrule
CICIoT2023 & 12 & 0 & - & 30 & 30 & 9.999 & 9.999 & 0.9935 & 0.9935 & +0.0000 & +0.0000 & v2\_\_ciciot\_\_sens\_chi0\_dp\_\_seed12 \\
CICIoT2023 & 13 & 1 & 21 & 29 & 30 & 9.858 & 9.999 & 0.9911 & 0.9912 & -0.0000 & -0.0000 & v2\_\_ciciot\_\_sens\_chi0\_dp\_\_seed13 \\
CICIoT2023 & 14 & 0 & - & 30 & 30 & 9.999 & 9.999 & 0.9939 & 0.9939 & +0.0000 & +0.0000 & v2\_\_ciciot\_\_sens\_chi0\_dp\_\_seed14 \\
CICIoT2023 & 15 & 4 & 14;15;16;22 & 26 & 30 & 9.434 & 9.999 & 0.9951 & 0.9946 & +0.0006 & +0.0000 & v2\_\_ciciot\_\_sens\_chi0\_dp\_\_seed15 \\
CICIoT2023 & 16 & 0 & - & 30 & 30 & 9.999 & 9.999 & 0.9919 & 0.9919 & +0.0000 & +0.0000 & v2\_\_ciciot\_\_sens\_chi0\_dp\_\_seed16 \\
UNSW-NB15 & 12 & 0 & - & 30 & 30 & 9.999 & 9.999 & 0.9569 & 0.9569 & +0.0000 & +0.0000 & v2\_\_unsw\_\_sens\_chi0\_dp\_\_seed12 \\
UNSW-NB15 & 13 & 1 & 21 & 29 & 30 & 9.858 & 9.999 & 0.9627 & 0.9603 & +0.0024 & +0.0018 & v2\_\_unsw\_\_sens\_chi0\_dp\_\_seed13 \\
UNSW-NB15 & 14 & 0 & - & 30 & 30 & 9.999 & 9.999 & 0.9389 & 0.9389 & +0.0000 & +0.0000 & v2\_\_unsw\_\_sens\_chi0\_dp\_\_seed14 \\
UNSW-NB15 & 15 & 4 & 14;15;16;22 & 26 & 30 & 9.434 & 9.999 & 0.9388 & 0.9322 & +0.0067 & +0.0034 & v2\_\_unsw\_\_sens\_chi0\_dp\_\_seed15 \\
UNSW-NB15 & 16 & 0 & - & 30 & 30 & 9.999 & 9.999 & 0.9211 & 0.9211 & +0.0000 & -0.0000 & v2\_\_unsw\_\_sens\_chi0\_dp\_\_seed16 \\
\bottomrule
\end{tabular}
\end{table*}


% Figure environments (labels used above)
\begin{figure}[!t]\centering\includegraphics[width=\linewidth]{figures/FIG1_attack_robustness.pdf}\caption{Attack robustness (dots: seeds; bars: means). Static non-adaptive attacks only.}\label{fig:supp-attack}\end{figure}
\begin{figure}[!t]\centering\includegraphics[width=\linewidth]{figures/FIG2_sensitivity_forest.pdf}\caption{Paired change in selected-checkpoint ROC--AUC versus the parent arm (95\% $t$-interval).}\label{fig:supp-sens}\end{figure}
\begin{figure}[!t]\centering\includegraphics[width=\linewidth]{figures/FIG3_routing_x_surrogate.pdf}\caption{Routing$\times$surrogate design.}\label{fig:supp-int}\end{figure}
\begin{figure}[!t]\centering\includegraphics[width=\linewidth]{figures/FIG4_participation_fairness.pdf}\caption{Participation: held-out client recall and selection Gini.}\label{fig:supp-part}\end{figure}
\begin{figure}[!t]\centering\includegraphics[width=\linewidth]{figures/FIG5_scalability.pdf}\caption{Scalability (3 seeds).}\label{fig:supp-scale}\end{figure}
